\section{Introduction}
\label{sec:intro}

Can we change a single fact inside a language model without breaking anything else?
Knowledge editing methods claim exactly this: surgical, localized updates that modify one behavior while preserving all others \citep{meng2022rome,meng2023memit,fang2025alphaedit}.
If this works, it would enable safe post-deployment corrections---fixing outdated facts, removing harmful associations, or patching errors---without expensive retraining.
But if it does not, then every ``fix'' risks unpredictable side effects, undermining the reliability of deployed models.

We probe this question with a deliberately simple test case: \textbf{can we teach \gptxl that $2{+}2{=}5$ without changing any other behavior?}
Arithmetic is an ideal stress test for edit isolation because it involves \emph{computational} knowledge---the model must execute an operation, not merely recall a stored association.
If editing methods can isolate even this simple change, their claims of surgical precision are strengthened.
If they cannot, it reveals fundamental limits in how knowledge is stored and modified in neural networks.

Existing knowledge editing benchmarks (\counterfact, zsRE, \knowedit) evaluate exclusively on \emph{factual} edits---changing entity-relation tuples such as ``The Eiffel Tower is in [Paris $\rightarrow$ London]'' \citep{meng2022rome,zhang2024comprehensive}.
Locality metrics in these benchmarks measure whether semantically \emph{unrelated} facts are preserved, but they do not test whether computationally \emph{related} operations are affected.
No prior work systematically evaluates whether editing arithmetic knowledge produces side effects on neighboring arithmetic operations.

We apply three editing methods to \gptxl: \rome \citep{meng2022rome}, fine-tuning (last 4 layers), and \lora \citep{biderman2024lora}.
We evaluate each on a custom arithmetic suite measuring (1) edit efficacy, (2) paraphrase generalization, (3) side effects on 15 related arithmetic queries, (4) side effects on 15 unrelated arithmetic queries, and (5) general language modeling via perplexity.
Our key finding is striking: \textbf{all three methods introduce a systematic \fivebias}, causing nearly every arithmetic prompt to output ``5'' after editing.
This bias is not random noise---it is a consistent, directional shift toward the injected answer, suggesting that arithmetic knowledge relies on shared output-digit circuits that resist isolated modification.

\rome achieves the best isolation on general language modeling (perplexity essentially unchanged at 10.22 vs.\ 10.30), but still corrupts arithmetic broadly.
\lora causes the most damage, increasing perplexity $4\times$ while propagating the bias to all tested paraphrases.
The uniformity of this failure across three fundamentally different editing approaches---rank-one weight update, gradient descent, and low-rank adaptation---suggests the limitation is architectural, not methodological.

We make the following contributions:
\begin{itemize}[leftmargin=*,itemsep=0pt,topsep=0pt]
    \item We conduct the first systematic study of editing \emph{arithmetic} knowledge in LLMs, revealing that computational knowledge behaves fundamentally differently from factual knowledge under editing.
    \item We identify a \emph{systematic} \fivebias phenomenon: all tested editing methods cause the injected digit to dominate arithmetic outputs, providing evidence that arithmetic uses shared computational circuits.
    \item We provide a fine-grained per-query analysis showing that side effects are predictable and directional, not random, with implications for understanding how transformers represent arithmetic.
\end{itemize}
